\setcounter{aufgabe}{16}
\hypertarget{e3}{}\einheitenkopf{Einheit 3 von 5 · Normalform und p-q-Formel}

\begin{aufgabe}{p-q-Formel -- Welche Form? Kreuze für jede Gleichung die Form an.
Löse nichts.}

\sachtabelle{lcccc}{Gleichung & nur $x^2$ & Produkt $=0$ & Klammer im Quadrat &
Normalform}{%
a) \ $x^2 = 121$ & $\square$ & $\square$ & $\square$ & $\square$ \\
b) \ $x^2 + 9x - 22 = 0$ & $\square$ & $\square$ & $\square$ & $\square$ \\
c) \ $(x - 6) \cdot (x + 1) = 0$ & $\square$ & $\square$ & $\square$ & $\square$ \\
d) \ $(x + 8)^2 = 25$ & $\square$ & $\square$ & $\square$ & $\square$ \\
e) \ $x^2 - 3x = 0$ & $\square$ & $\square$ & $\square$ & $\square$}
\end{aufgabe}

\begin{aufgabe}{p-q-Formel -- Was sind $p$ und $q$? Lies beide Zahlen mit Vorzeichen
ab. Rechne nichts aus.}
\begin{geruest}
\gz{a}{$x^2 + 10x + 3 = 0$}{\feld{p}\feld{q}}
\gz{b}{$x^2 - 4x + 9 = 0$}{\feld{p}\feld{q}}
\gz{c}{$x^2 + 7x - 5 = 0$}{\feld{p}\feld{q}}
\gz{d}{$x^2 - 9x - 2 = 0$}{\feld{p}\feld{q}}
\gz{e}{$x^2 - x + 6 = 0$}{\feld{p}\feld{q}}
\end{geruest}
Bring die Gleichung erst auf die Form $\ldots = 0$, dann lies ab.
\begin{geruest}
\gz{f}{$x^2 + 15 = 8x$}{\feld{p}\feld{q}}
\end{geruest}
\end{aufgabe}

\begin{aufgabe}{p-q-Formel -- Vorzahl von $x^2$? Kreuze an, ob vor der Lösungsformel
geteilt werden muss, und trag die Zahl ein. Rechne nichts.}
\begin{geruest}
\gz{a}{$x^2 + 5x - 6 = 0$}{\janein\ durch \leerfeld}
\gz{b}{$3x^2 - 9x + 6 = 0$}{\janein\ durch \leerfeld}
\gz{c}{$x^2 - 11x = 0$}{\janein\ durch \leerfeld}
\gz{d}{$-x^2 + 4x + 5 = 0$}{\janein\ durch \leerfeld}
\gz{e}{$5x^2 + 20x - 25 = 0$}{\janein\ durch \leerfeld}
\end{geruest}
\end{aufgabe}

\begin{aufgabe}{p-q-Formel -- rechnen. Lies $p$ und $q$ mit Vorzeichen ab und setze in
die Lösungsformel der Formelsammlung ein. Gib beide Lösungen an.}
\beispiel{x^2 + 6x - 16 &= 0 &&\mid p = 6,\ q = -16 \\
x_{1,2} &= -3 \pm \sqrt{9 + 16} \\
x_{1,2} &= -3 \pm 5 \\
x_1 &= 2 && x_2 = -8}
\begin{gleichungsraster}[3]
\gl{x^2 + 8x + 15 = 0} & \gl{x^2 + 10x + 21 = 0} \\
\gl{x^2 + 6x + 5 = 0} & \gl{x^2 + 12x + 32 = 0} \\
\gl{x^2 + 4x - 21 = 0} & \gl{x^2 - 10x + 24 = 0} \\
\gl{x^2 - 2x - 35 = 0} & \gl{x^2 + 5x + 6 = 0} \\
\end{gleichungsraster}
\end{aufgabe}

\begin{aufgabe}{p-q-Formel -- wie viele Lösungen? Rechne zuerst den Wert unter der
Wurzel aus und schreib hin, wie viele Lösungen es gibt. Rechne die Lösungen danach
aus, soweit es welche gibt. Geht die Wurzel nicht auf, lass sie stehen und gib
zusätzlich den Näherungswert auf zwei Stellen an.}
\begin{gleichungsraster}[4]
\gl{x^2 - 14x + 49 = 0} & \gl{x^2 + 4x + 9 = 0} \\
\gl{x^2 - 8x + 11 = 0} & \\
\end{gleichungsraster}
\end{aufgabe}

\begin{aufgabe}{p-q-Formel -- ordnen. Bring erst alles auf eine Seite und ordne nach
$x^2$, $x$ und Zahl. Löse dann.}
\begin{gleichungsraster}[4]
\gl{x^2 + 5x = 24} & \gl{x^2 = 7x + 18} \\
\end{gleichungsraster}
Löse zuerst die Klammer auf, dann ordne.
\begin{gleichungsraster}[5]
\gl{(x + 3)^2 = 4x + 21} & \\
\end{gleichungsraster}
\end{aufgabe}

\begin{aufgabe}{p-q-Formel -- normieren. Teile zuerst durch die Vorzahl von $x^2$,
dann löse.}
\beispiel{2x^2 - 6x - 20 &= 0 &&\mid : 2 \\
x^2 - 3x - 10 &= 0 \\
x_{1,2} &= 1{,}5 \pm \sqrt{2{,}25 + 10} = 1{,}5 \pm 3{,}5 \\
x_1 &= 5 && x_2 = -2}
\begin{gleichungsraster}[4]
\gl{2x^2 + 14x + 12 = 0} & \gl{3x^2 - 18x + 24 = 0} \\
\gl{-x^2 + 2x + 8 = 0} & \\
\end{gleichungsraster}
\end{aufgabe}

\begin{aufgabe}{p-q-Formel -- Probe. Löse die Gleichung und setze danach eine deiner
Lösungen zur Probe wieder ein.}
\begin{gleichungsraster}[5]
\gl{x^2 - 9x + 20 = 0} & \gl{x^2 + 3x - 10 = 0} \\
\end{gleichungsraster}
\end{aufgabe}

\begin{aufgabe}{p-q-Formel -- aus der Prüfung. Überleg zuerst, was nötig ist: ordnen,
teilen oder beides. Löse dann.}
\begin{gleichungsraster}[5]
\gl{2x^2 + 14x + 24 = 0} & \gl{-x^2 - 4x + 7 = -5} \\
\end{gleichungsraster}
\hfill (P10 2020) \qquad (P10 2023)

Gib die Lösungen zuerst mit Wurzel und dann als Näherungswert auf zwei Stellen an.
\begin{gleichungsraster}[5]
\gl{x^2 - 10x + 23 = 0} & \\
\end{gleichungsraster}
\hfill (P10 2025)
\end{aufgabe}

\begin{aufgabe}{p-q-Formel -- Fehler finden. Tom löst so:}
\rechnung{x^2 - 8x + 12 &= 0 &&\mid p = -8,\ q = 12 \\
x_{1,2} &= -4 \pm \sqrt{16 - 12} \\
x_1 &= -2 && x_2 = -6}
Finde den Fehler, benenne ihn und korrigiere die Rechnung.

\begin{teile}
\teil Löse $x^2 - 12x + 27 = 0$.
\end{teile}
\end{aufgabe}

\begin{aufgabe}{p-q-Formel -- Begründen. Antworte ohne zu rechnen.}
\begin{teile}
\teil Begründe, warum die Lösungsformel erst benutzt werden darf, wenn vor dem $x^2$
eine Eins steht.
\teil Bei $x^2 + 6x + 11 = 0$ ist der Wert unter der Wurzel negativ. Was bedeutet das
für die Lösungen? Begründe.
\teil Kann eine quadratische Gleichung genau drei Lösungen haben? Begründe mit dem
Aufbau der Lösungsformel.
\end{teile}
\end{aufgabe}

\begin{aufgabe}{p-q-Formel -- Lösungsweg wählen. Kreuze an, womit du am schnellsten
zum Ziel kommst. Rechne nicht.}
\begin{teile}
\teil $x^2 = 144$ \hfill \kreuz{Wurzelziehen}\kreuz{Nullprodukt}\kreuz{Lösungsformel}
\teil $x^2 - 15x = 0$ \hfill \kreuz{Wurzelziehen}\kreuz{Nullprodukt}\kreuz{Lösungsformel}
\teil $x^2 + 9x + 14 = 0$ \hfill \kreuz{Wurzelziehen}\kreuz{Nullprodukt}\kreuz{Lösungsformel}
\teil $(x - 4)^2 = 49$ \hfill \kreuz{Wurzelziehen}\kreuz{Nullprodukt}\kreuz{Lösungsformel}
\teil $(x + 2) \cdot (x - 11) = 0$ \hfill \kreuz{Wurzelziehen}\kreuz{Nullprodukt}\kreuz{Lösungsformel}
\end{teile}
\end{aufgabe}
